% Imporditud: tools/import_eksperimendid.py
% Allikas: Eesti füüsikaolümpiaadi eksperimendi ülesannete
%   kogu (Rael Kalda, 2023), ülesanne Ü3, lahendus L3.
\setAuthor{EFO žürii}
\setRound{lõppvoor}
\setYear{1997}
\setNumber{P E1}
\setDifficulty{1}
\setTopic{vedelike mehaanika}
\setVahendid{tops, joonlaud}
\setPunktid{8}
\setAllikas{Eksperimendi kogu Ü3, lahendus lk 34}
\setLahendusseis{toores}

\prob{Topsi rõhk}
Kui suurt rõhku avaldab lauale veega täidetud tops?

\hint

\solu
Teooria kirjeldus (2 p) Teoreetiline pōhjendus seisneb vōrrandite süsteemi lahendamises, mis saadakse Hookei seaduse, hōōrdejōu seaduse ja raskusjōu valemi rakendamise tulemusena.

Katse kirjeldus ja mōōtmised (2 p)

Mōōtevea hindamine (2. p)

Järeldused (2 p)
\probend
